% !TeX root = ../../Book4.tex
% 纲要标题：### A.3 Barr–Beck 定理的假设与典型例子
% 纲要位置：工作记录/计划纲要.md，第 3836 行。
% 纲要细目（写作范围）：
% #### A.3.1 分裂余等化子
% #### A.3.2 Barr–Beck 单子性判据
% #### A.3.3 模的限制标量与单子性
% #### A.3.4 忠实平坦下降的含义


\section{Barr–Beck 定理的假设与典型例子}
\label{b4:sec:A03}

比较函子总能构造，但它何时为等价，需要一个可以检验的条件。基本观察是：每个单子代数都能由两个自由代数之间的余等化子给出，而且遗忘结构后，这个余等化子具有显式分裂。下面只要求原范畴正确处理这种特殊的余等化子。

\subsection{分裂余等化子}
\label{b4:subsec:A03:01}

\begin{definition}\label{b4:def:A-split-coequalizer}
在范畴 \(\mathcal C\) 中，设
\[
A\mathrel{\substack{\xrightarrow{\ f\ }\\[-0.4ex]\xrightarrow[\ g\ ]{}}}B
\xrightarrow{q}C,\qquad qf=qg.
\]
如果还存在 \(s:C\to B\)、\(t:B\to A\)，满足
\begin{equation}\label{b4:eq:A-split-coequalizer}
qs=1_C,\qquad ft=1_B,\qquad gt=sq,
\end{equation}
则称这个叉形图连同 \(s,t\) 为一个\term[fenlie-yudenghuazi]{分裂余等化子}[split coequalizer]。
\end{definition}

等式 \(qs=1_C\)、\(ft=1_B\) 分别给出 \(q,f\) 的截面；还必须满足 \(gt=sq\)，使沿 \(t\) 返回 \(A\) 后再施行 \(g\)，与先施行 \(q\) 再沿 \(s\) 返回相同。若两张平行箭头只有共同截面 \(r:B\rightarrow A\)，已知的只是 \(fr=gr=1_B\)。这与 \(gt=sq\) 所要求的相容性不同，不能用共同截面代替定义中的全部等式。

\begin{lemma}\label{b4:lem:A-split-absolute}
设 \(\mathcal C\) 为范畴，\(f,g:A\rightrightarrows B\)、\(q:B\to C\) 为其中的态射。若 \(qf=qg\)，并存在 \(s:C\to B\)、\(t:B\to A\) 满足 \(qs=1_C\)、\(ft=1_B\)、\(gt=sq\)，则 \(q\) 确实是 \(f,g\) 的余等化子。任意函子 \(H:\mathcal C\to\mathcal E\)（\(\mathcal E\) 为任意范畴）都保持这个余等化子及其分裂。
\end{lemma}
\begin{proof}
若 \(h:B\to Z\) 满足 \(hf=hg\)，则
\[
h=hft=hgt=hsq.
\]
所以 \(hs:C\to Z\) 给出所需分解。若 \(uq=h\)，则 \(u=uqs=hs\)，因而唯一。任意函子都保持复合、恒等及所列等式，对像图应用同一证明即可。
\end{proof}

对函子 \(U:\mathcal D\to\mathcal C\)，如果平行态射 \(f,g:A\rightrightarrows B\) 的像在 \(\mathcal C\) 中可以补成一个分裂余等化子，就称它们是 \(U\)-分裂的。分裂态射只要求存在于 \(\mathcal C\)，不要求来自 \(\mathcal D\)。

\begin{proposition}\label{b4:prop:A-em-split-coequalizers}
设 \((T,\eta,\mu)\) 为局部小范畴 \(\mathcal C\) 上的单子，\(\mathcal C^T\) 为其代数范畴，\(U^T:\mathcal C^T\to\mathcal C\) 为遗忘函子。若 \(\mathcal C^T\) 中一对平行代数同态
\[
f,g:(A,a)\rightrightarrows(B,b)
\]
的底层图在 \(\mathcal C\) 中有分裂余等化子 \(q:B\to C\)，则 \(C\) 上存在唯一代数结构 \(c:TC\to C\)，使 \(q\) 为代数同态；带此结构的 \(q\) 是 \(\mathcal C^T\) 中的余等化子。此外，\(U^T\) 反映同构。
\end{proposition}
\begin{proof}
由分裂等式 \(qs=1_C\)，施行 \(T\) 及 \(T^2\) 得到
\[
Tq\,Ts=1_{TC},\qquad T^2q\,T^2s=1_{T^2C}.
\]
所以 \(q,Tq,T^2q\) 都是分裂满态射；这里不要求 \(T\) 保持一般满态射。\cref{b4:lem:A-split-absolute}又说明 \(Tq\) 是 \(Tf,Tg\) 的余等化子。于是
\[
qbT(f)=qfa=qga=qbT(g),
\]
所以存在唯一 \(c:TC\to C\)，使 \(cTq=qb\)。这正是 \(q\) 成为代数同态的条件。

由于 \(q\) 为满态射，验证 \(c\eta_C=1_C\) 可以在右边复合 \(q\)；此时由单位自然性得到
\[
c\eta_Cq=cTq\eta_B=qb\eta_B=q.
\]
结合律则在右边复合 \(T^2q\)。这个态射也是分裂满态射，且
\[
cT(c)T^2q=qbT(b),\qquad
c\mu_CT^2q=qb\mu_B.
\]
两者由 \(b\) 的代数公理相等，故 \(cT(c)=c\mu_C\)。

若代数同态 \(h:(B,b)\to(D,d)\) 满足 \(hf=hg\)，则底层余等化子的泛性质给出唯一 \(v:C\to D\)，使 \(vq=h\)。复合满态射 \(Tq\) 后，
\[
vcTq=vqb=hb=dT(h)=dT(v)Tq,
\]
所以 \(vc=dT(v)\)，即 \(v\) 仍是代数同态。这证明了余等化子的泛性质。

最后，若代数同态 \(u:(X,a)\to(Y,b)\) 的底层态射可逆，从 \(ua=bT(u)\) 两边乘以逆，得到 \(u^{-1}b=aT(u^{-1})\)。因此底层逆也是代数同态，\(U^T\) 反映同构。
\end{proof}

\subsection{Barr–Beck 单子性判据}
\label{b4:subsec:A03:02}

如果函子把某个态射送成同构就能推出该态射本身为同构，称这个函子\term[baoshou-hanzi]{保守}[conservative]，也称它反映同构。这个条件不同于忠实：忠实只要求每个同态集合上的映射为单射。

\begin{theorem}[Barr–Beck 单子性判据]\label{b4:thm:A-beck-monadicity}
设 \(\mathcal C,\mathcal D\) 为局部小范畴，\(U:\mathcal D\to\mathcal C\) 有左伴随 \(F:\mathcal C\to\mathcal D\)，伴随的单位、余单位为 \(\eta,\varepsilon\)，所生单子为 \(T=UF\)。则比较函子 \(K:\mathcal D\to\mathcal C^T\) 全忠实且本质满，当且仅当下列条件均成立：
\begin{enumerate}
\item \(U\) 反映同构；
\item 每一对 \(U\)-分裂平行态射在 \(\mathcal D\) 中都有余等化子，且 \(U\) 保持该余等化子。
\end{enumerate}
若这些条件成立，并且已为每个 \(T\)-代数 \((X,a)\) 同时指定平行对
\(F(a),\varepsilon_{FX}:FTX\rightrightarrows FX\) 的余等化子
\(q_{(X,a)}:FX\to L_{(X,a)}\)，则这些已选数据给出 \(K\) 的拟逆，因而 \(K\) 是范畴等价。反过来，若已给出 \(K\) 的拟逆，这族余等化子也可由该拟逆同时构造。
\end{theorem}
这是 \term[barr-beck-danzixing]{Barr–Beck 单子性定理}[Barr--Beck monadicity theorem]的分裂余等化子版本\footnote{巴尔（Michael Barr，1937-），美国数学家，研究范畴论与同调代数；贝克（Jonathan Mock Beck，1935-2006），美国数学家，建立以其命名的单子性判据。}。第二项只涉及指定的一类平行态射，而不是要求原范畴存在并保持全部余等化子。
\begin{proof}
先设 \(K\) 全忠实且本质满。由\cref{b4:prop:A-em-split-coequalizers}，\(U^T\) 反映同构；全忠实的 \(K\) 也反映同构，而 \(U=U^TK\)，故 \(U\) 反映同构。

给定 \(\mathcal D\) 中的 \(U\)-分裂平行对 \(f,g:A\rightrightarrows B\)，并在 \(\mathcal C\) 中取其分裂余等化子 \(q:UB\to C\)。同一命题在 \(C\) 上给出代数结构 \(c\)，并使
\(\widehat q:K(B)\to(C,c)\) 成为 \(K(f),K(g)\) 的余等化子，其底层态射就是 \(q\)。本质满性给出对象 \(L\in\mathcal D\) 及同构
\(\psi:K(L)\xrightarrow{\sim}(C,c)\)；全忠实性给出唯一 \(h:B\to L\)，满足
\(K(h)=\psi^{-1}\widehat q\)。它余等化 \(f,g\)。对任何余等化它们的 \(r:B\to Z\)，\(K(r)\) 经 \(\widehat q\) 唯一分解；再复合 \(\psi\)，利用 \(K\) 的全忠实性，得到 \(r\) 经 \(h\) 的唯一分解。因此 \(h\) 是 \(\mathcal D\) 中的余等化子。又
\(q=U^T(\psi)U(h)\)，且 \(U^T(\psi)\) 为同构，故 \(U(h)\) 也是底层余等化子。其他余等化子与 \(h\) 唯一同构，所以也被 \(U\) 保持。这证明第二项；特别地，当 \(K\) 已是等价时，这一构造正是沿等价运输代数范畴中由 \(U^T\) 创建的分裂余等化子。

反过来，假定两项条件成立。先证明每个 \(T\)-代数都由比较函子得到。对 \((X,a)\)，代数结合律比较的是 \(T^2X\) 上的两种求值次序：\(\mu_X\) 先归并两层形式表达式，\(T(a)\) 先对内层求值，最后都再施行 \(a\)。在有限表单子的例子中，它们把 \([[x_1,x_2],[x_3]]\) 分别送到
\[
[x_1,x_2,x_3],\qquad
[a([x_1,x_2]),a([x_3])],
\]
而 \(a\mu_X=aT(a)\) 要求这两张表的最终值相同。因此在自由对象上应当识别的，正是下面两支平行箭头：
\begin{equation}\label{b4:eq:A-free-presentation}
FTX
\mathrel{\substack{\xrightarrow{\ \varepsilon_{FX}\ }\\[-0.4ex]
\xrightarrow[\ F(a)\ ]{}}}
FX.
\end{equation}
应用 \(U\) 后，它们是 \(\mu_X,T(a):T^2X\rightrightarrows TX\)。代数公理说明 \(a:TX\to X\) 余等化这对态射。取 \(s=\eta_X\)、\(t=\eta_{TX}\)，便有
\[
a\eta_X=1_X,\qquad
\mu_X\eta_{TX}=1_{TX},\qquad
T(a)\eta_{TX}=\eta_Xa.
\]
因此这是一个分裂余等化子。由第二项假设，\eqref{b4:eq:A-free-presentation}存在余等化子 \(q:FX\to L\)，且 \(Uq:TX\to UL\) 也是底层余等化子。余等化子的唯一性给出唯一同构
\(\phi:UL\xrightarrow{\sim}X\)，满足 \(\phi Uq=a\)。记
\(c_L=U(\varepsilon_L):T(UL)\rightarrow UL\)。此时只恢复了正确的底层对象，还须确认 \(K(L)\) 的求值结构也与 \(a\) 相同。把其结构经 \(\phi\) 运输到 \(X\)，得到
\[
b=\phi c_LT(\phi^{-1}):TX\longrightarrow X.
\]
余单位关于 \(q\) 的自然性为
\(q\varepsilon_{FX}=\varepsilon_LF(Uq)\)。应用 \(U\) 并在左侧复合 \(\phi\)，得到两个从 \(T^2X\) 到 \(X\) 的态射相等：
\[
 a\mu_X=\phi c_LT(Uq)=bT(a).
\]
代数公理又给出 \(a\mu_X=aT(a)\)，故 \(bT(a)=aT(a)\)。由单位律，
\[
T(a)T(\eta_X)=T(a\eta_X)=1_{TX}.
\]
在两边右复合 \(T(\eta_X)\)，便得 \(b=a\)。于是
\(\phi c_L=aT(\phi)\)，所以 \(\phi\) 是代数同构
\(K(L)\xrightarrow{\sim}(X,a)\)，证明本质满性。

要从底层代数同态恢复 \(\mathcal D\) 中的态射，需要先证明：每个 \(Y\in\mathcal D\) 的余单位 \(\varepsilon_Y\) 就是其自由展示的余等化子。写 \(a_Y=U(\varepsilon_Y)\)。相同构造给出
\[
FTUY\rightrightarrows FUY\xrightarrow{q}Y',
\]
其中两箭头为 \(\varepsilon_{FUY}\)、\(F(a_Y)\)，而 \(Uq\) 是余等化子。余单位自然性说明 \(\varepsilon_Y:FUY\to Y\) 也余等化这对态射，故唯一分解为 \(hq\)。\(U(\varepsilon_Y)=a_Y\) 本身是上面已经验证的分裂余等化子，所以 \(Uh\) 为两个余等化子之间的同构。由第一项假设，\(h\) 为同构。因此 \(\varepsilon_Y\) 就是这一对态射的余等化子，特别是满态射。

若 \(r,s:Y\to Z\) 满足 \(Ur=Us\)，则余单位自然性给出
\[
r\varepsilon_Y=\varepsilon_ZF(Ur)
=\varepsilon_ZF(Us)=s\varepsilon_Y.
\]
消去满态射 \(\varepsilon_Y\)，得到 \(r=s\)，故 \(U\) 以及 \(K\) 忠实。

现在给定代数同态 \(\alpha:K(Y)\to K(Z)\)，即
\(\alpha a_Y=a_ZT(\alpha)\)。令 \(\widehat\alpha=\varepsilon_ZF(\alpha):FUY\to Z\)。两次使用余单位自然性可得
\[
\begin{aligned}
\widehat\alpha\,\varepsilon_{FUY}
&=\varepsilon_Z\varepsilon_{FUZ}FT(\alpha)\\
&=\varepsilon_ZF(a_Z)FT(\alpha)\\
&=\varepsilon_ZF(\alpha a_Y)
=\widehat\alpha\,F(a_Y).
\end{aligned}
\]
因此它经过余等化子 \(\varepsilon_Y\) 唯一分解为 \(r\varepsilon_Y\)，其中 \(r:Y\to Z\)。应用 \(U\)，得到
\[
(Ur)a_Y=a_ZT(\alpha)=\alpha a_Y.
\]
\(a_Y\) 有右逆 \(\eta_{UY}\)，消去后便有 \(Ur=\alpha\)。这证明全性。

若自由展示的余等化子 \(q_{(X,a)}\) 已同时指定，刚才的运输给出唯一的代数同构
\(K(L_{(X,a)})\xrightarrow{\sim}(X,a)\)。\cref{b4:thm:equivalence-criterion}中带有已选代表的判据便给出拟逆；它在态射上由 \(K\) 的全忠实性唯一确定。对于大范畴，所需代表由规定的余等化子构造或上述已选数据给出。

反过来，若已有拟逆 \(J:\mathcal C^T\to\mathcal D\) 及自然同构
\(\alpha:1_{\mathcal D}\xrightarrow{\sim}JK\)，对代数 \((X,a)\) 取
\[
q_{(X,a)}=J(a)\alpha_{FX}:FX\longrightarrow J(X,a).
\]
这里 \(a:F^T(X)\to(X,a)\) 是自由代数展示的余等化子，而 \(KF=F^T\)。拟逆保持余等化子，\(\alpha\) 的自然性又识别两条平行箭头，故这些公式同时给出所需余等化子。
\end{proof}

离散拓扑伴随的反例现在有了明确解释：底集合函子不反映同构，连续双射可能不是同胚。对其他例子，即使已核验保守性，也仍须检查第二项指定的余等化子条件。

\subsection{模的限制标量与单子性}
\label{b4:subsec:A03:03}

\begin{proposition}\label{b4:prop:A-restriction-scalars-monadic}
设 \(R\to S\) 为含幺环同态，以下模范畴均以幺性左模及模同态为对象和态射。限制标量函子
\[
U:S\text{-}\mathbf{Mod}\longrightarrow R\text{-}\mathbf{Mod}
\]
具有单子性。相应单子在幺性左 \(R\) 模 \(M\) 上为 \(T(M)=S\otimes_RM\)，其中 \(S\) 的右 \(R\) 作用与所得张量积的左 \(R\) 作用均由 \(R\to S\) 给出；单位为 \(m\mapsto1\otimes m\)，乘法为 \(s\otimes t\otimes m\mapsto st\otimes m\)。
\end{proposition}
\begin{proof}
左伴随为扩张标量 \(F(M)=S\otimes_RM\)：映射
\[
\operatorname{Hom}_S(S\otimes_RM,N)\longrightarrow
\operatorname{Hom}_R(M,UN),\qquad
f\longmapsto\bigl(m\mapsto f(1\otimes m)\bigr)
\]
的逆为 \(h\mapsto\bigl(s\otimes m\mapsto s\cdot h(m)\bigr)\)。平衡关系及 \(R\)-线性保证逆映射良定义。

一个 \(S\)-线性映射若在底层为可逆 \(R\)-线性映射，其逆仍保持 \(S\)-作用，所以 \(U\) 反映同构。对 \(S\)-线性平行态射 \(f,g:M\rightrightarrows N\)，余等化子为
\[
N/\operatorname{im}(f-g).
\]
像本来就是 \(S\)-子模；遗忘到 \(R\) 后，底层像与商均不变，仍是 \(R\)-模中的余等化子。因此全部余等化子都存在并被 \(U\) 保持。上述规定的商模构造同时给出了自由展示的余等化子，故\cref{b4:thm:A-beck-monadicity}给出范畴等价。
\end{proof}

这里的代数结构也能直接辨认。对 \(R\)-模 \(M\)，一个 \(T\)-代数态射 \(a:S\otimes_RM\to M\) 定义 \(s\cdot m=a(s\otimes m)\)。单位律给出 \(1m=m\)，结合律给出 \((st)m=s(tm)\)，平衡关系使它延伸原来的 \(R\)-作用。反向由任何 \(S\)-模作用都得到这样的 \(a\)。这个直接说明与比较函子恰好一致。

\subsection{忠实平坦下降的含义}
\label{b4:subsec:A03:04}

限制标量的单子性对任意环同态成立。下降研究另一个方向：从已经扩张标量的对象以及相容数据，恢复原来的对象。因此这里使用扩张标量函子的余单子，而不是把上一命题再改写一次。

\begin{definition}\label{b4:def:A-faithfully-flat-descent-data}
设 \(R\to S\) 为交换环同态，\(M\) 为 \(S\)-模。若 \(S\)-线性映射
\(\gamma:M\to S\otimes_RM\) 满足
\begin{equation}\label{b4:eq:A-module-descent-data}
\begin{aligned}
e\gamma&=1_M,\qquad e(s\otimes m)=sm,\\
(1_S\otimes\gamma)\gamma&=\delta\gamma,\qquad
\delta(s\otimes m)=s\otimes1\otimes m,
\end{aligned}
\end{equation}
则称 \((M,\gamma)\) 为一个以余单子形式给出的\term[mo-xiajiang-shuju]{模下降数据}[module descent datum]。这里 \(S\) 在张量积的第一个因子上作用。对两份数据 \((M,\gamma)\)、\((M',\gamma')\)，若 \(S\) 线性映射 \(f:M\to M'\) 满足 \((1_S\otimes f)\gamma=\gamma'f\)，则称 \(f\) 为它们之间的同态。
\end{definition}

若 \(S=R\times R\)，一个 \(S\)-模分成两块 \(M_1\oplus M_2\)，而从 \(R\)-模 \(V\) 扩张得到的是 \(V\oplus V\)。因此要判断给定的两块是否来自同一个模，还须指定它们怎样互相识别。\cref{b4:prop:A-product-descent-model}将证明，\(\gamma\) 恰好记录一个同构 \(M_1\xrightarrow{\sim}M_2\)；余结合律使两方向的识别互逆。

\ProseParagraphBreak

这些条件正是伴随 \(S\otimes_R-\dashv U\) 所生余单子的余代数公理。对 \(R\)-模 \(V\)，扩张模 \(S\otimes_RV\) 总带有这样的数据：
\[
\gamma_V(s\otimes v)=s\otimes(1\otimes v).
\]
等式逐个在纯张量上成立。问题是给定的数据是否都这样产生，以及保持数据的态射是否也来自唯一的 \(R\)-线性映射。

\begin{theorem}[忠实平坦模下降]\label{b4:thm:A-faithfully-flat-descent}
设 \(R\to S\) 为交换含幺环同态，且 \(S\) 是忠实平坦 \(R\) 模，即 \(S\otimes_R-\) 正合，并对每个 \(R\) 模 \(V\) 满足
\(S\otimes_RV=0\Rightarrow V=0\)。以 \(\mathsf{Desc}(S/R)\) 表示 \(S/R\) 模下降数据及其同态组成的范畴。令 \(\gamma_V(s\otimes v)=s\otimes(1\otimes v)\)，则
\[
V\longmapsto(S\otimes_RV,\gamma_V)
\]
是 \(R\) 模范畴到 \(\mathsf{Desc}(S/R)\) 的等价，态射 \(f:V\to V'\) 被送到 \(1_S\otimes f\)。给定下降数据 \((M,\gamma)\)，一个恢复对象为在 \(R\)-模范畴中取的等化子
\begin{equation}\label{b4:eq:A-descent-equalizer}
V=\bigl\{m\in M:\gamma(m)=1\otimes m\bigr\},
\end{equation}
相应同构 \(S\otimes_RV\to M\) 为 \(s\otimes v\mapsto sv\)。
\end{theorem}
\symidx{\(\mathsf{Desc}(S/R)\)：模下降数据组成的范畴}
\begin{proof}
把\cref{b4:thm:A-beck-monadicity}应用于对偶范畴：左伴随 \(F=S\otimes_R-\) 具有余单子性，只须它反映同构，并保持所有在 \(F\) 下分裂的等化子。平坦性使 \(F\) 正合，从而保持模范畴中的全部等化子。若 \(F(f)\) 可逆，则正合性给出
\[
F(\ker f)=0,\qquad F(\operatorname{coker}f)=0.
\]
忠实性假设使原来的核、余核都为零，故 \(f\) 可逆。模范畴的等化子可统一取为两态射之差的核子模，给出对偶判据所需的已选对象。因此该判据得到范畴等价。

为说明所列恢复公式，记 \(\eta_M(m)=1\otimes m\)。\(\gamma\) 与 \(\eta_M\) 都是 \(R\)-线性的，但 \(\eta_M\) 一般不是 \(S\)-线性的：
\[
\eta_M(sm)=1\otimes sm,\qquad
s\cdot\eta_M(m)=s\otimes m,
\]
这两个张量未必相同。因此\eqref{b4:eq:A-descent-equalizer}首先是 \(R\)-模中的等化子，给恢复对象 \(V\) 以 \(R\)-模结构。平坦性使扩张后的 \(S\otimes_RV\) 成为 \(S\)-模中的等化子：
\[
1_S\otimes\gamma,\ 1_S\otimes\eta_M:
S\otimes_RM\rightrightarrows S\otimes_RS\otimes_RM
\]
两支箭头都保持第一个 \(S\) 因子上的作用。另一方面，\eqref{b4:eq:A-module-descent-data}说明 \(\gamma\) 的像属于这个等化子，且 \(e\gamma=1_M\)。

若 \(z\in S\otimes_RM\) 满足两映射取值相同，对该等式施行
\(s\otimes t\otimes m\mapsto st\otimes m\)。由于 \(\gamma\) 是 \(S\)-线性的，左边变成 \(\gamma\bigl(e(z)\bigr)\)，右边变成 \(z\)。因此 \(z=\gamma\bigl(e(z)\bigr)\)，说明 \(\gamma\) 把 \(M\) 同构到该等化子。将它与 \(S\otimes_RV\) 的等化子识别相比较，逆映射为
\(\phi:S\otimes_RV\to M\)、\(\phi(s\otimes v)=sv\)。它还保持下降数据：对 \(v\in V\)，
\[
\gamma\phi(s\otimes v)=\gamma(sv)=s\otimes v
=(1_S\otimes\phi)\gamma_V(s\otimes v).
\]
因此 \(\phi\) 是下降数据之间的同构。

对保持下降数据的态射 \(f:(M,\gamma)\to(M',\gamma')\)，若 \(v\in V\)，则
\(\gamma'f(v)=(1_S\otimes f)\gamma(v)=1\otimes f(v)\)，故 \(f\) 限制为恢复子模之间的 \(R\)-线性映射 \(f|_V:V\to V'\)。限制保持恒等与复合，而且
\(\phi'(1_S\otimes f|_V)=f\phi\)，所以所给恢复公式及其同构对态射自然。

最后，对原来的 \(R\)-模 \(W\)，令 \(V_W\) 为标准下降数据 \((S\otimes_RW,\gamma_W)\) 的恢复子模。映射 \(\theta_W:W\rightarrow V_W\)、\(w\mapsto1\otimes w\) 扩张标量后，与刚才的乘法同构复合为 \(1_{S\otimes_RW}\)，故 \(S\otimes_R\theta_W\) 可逆。前面的等化子识别用平坦性得到了 \(S\otimes_RV_W\cong S\otimes_RW\)；此处还要用忠实性，才能从扩张后核、余核消失推出 \(\theta_W\) 本身可逆。这给出另一方向的自然同构，因而上述恢复函子就是所列等价的拟逆。
\end{proof}

对角映射 \(R\rightarrow R\times R\) 是忠实平坦的，因为张量函子把 \(V\) 送到 \(V\oplus V\)，既保持正合列，也不会把非零模变为零。下面把这个例子的下降条件直接写成两块之间的映射。

\begin{proposition}\label{b4:prop:A-product-descent-model}
设 \(R\) 为交换含幺环，\(S=R\times R\)，\(R\rightarrow S\) 为对角映射。令 \(e_1=(1,0)\)、\(e_2=(0,1)\)。对幺性 \(S\)-模 \(M\)，写 \(M=M_1\oplus M_2\)，其中 \(M_i=e_iM\)。在同构 \(S\otimes_RM\cong M\oplus M\) 下，\(M\) 上的模下降数据与 \(R\)-模同构 \(f:M_1\rightarrow M_2\) 一一对应，相应结构为
\[
\gamma(m_1,m_2)
=\Bigl(\bigl(m_1,f(m_1)\bigr),
\bigl(f^{-1}(m_2),m_2\bigr)\Bigr).
\]
其恢复模为 \(f\) 的图
\[
V=\{\,(m_1,f(m_1))\mid m_1\in M_1\,\},
\]
以坐标投影同构于 \(M_1\)。
\end{proposition}
\begin{proof}
同构 \(S\otimes_RM\cong M\oplus M\) 把 \(e_1\otimes m\)、\(e_2\otimes m\) 分别送到 \((m,0)\)、\((0,m)\)；外层两块承受 \(S\) 的两个坐标作用。\(\gamma\) 的 \(S\)-线性因此给出 \(R\)-线性映射 \(\gamma_{ij}:M_j\rightarrow M_i\)，满足
\[
\gamma(m_1,m_2)
=\bigl((\gamma_{11}(m_1),\gamma_{21}(m_1)),
(\gamma_{12}(m_2),\gamma_{22}(m_2))\bigr).
\]
余单位 \(e\) 从外层第一块取内部 \(M_1\) 分量，从第二块取 \(M_2\) 分量，所以 \(e\gamma=1_M\) 等价于 \(\gamma_{11}=1_{M_1}\)、\(\gamma_{22}=1_{M_2}\)。

对 \(m\in M_j\)，比较余结合等式中外层 \(e_j\)、中间层 \(e_i\) 及内部 \(M_k\) 的分量：左边先用 \(\gamma_{ij}\)，再用 \(\gamma_{ki}\)；右边由 \(\delta(e_j\otimes m)=e_j\otimes1\otimes m\) 把中间层同时展开为 \(e_1+e_2\)。因此余结合律恰好要求
\[
\gamma_{ki}\gamma_{ij}=\gamma_{kj}
\qquad(i,j,k\in\{1,2\}).
\]
在两个对角映射已为恒等的条件下，唯一尚需验证的等式是
\[
\gamma_{12}\gamma_{21}=1_{M_1},\qquad
\gamma_{21}\gamma_{12}=1_{M_2}.
\]
故 \(f=\gamma_{21}\) 必为同构，且 \(\gamma_{12}=f^{-1}\)。反过来，任取这样的同构，所列公式是 \(S\)-线性的，并满足刚才全部分量等式，因而给出下降数据。

最后，\(1\otimes(m_1,m_2)\) 在外层两块中都是 \((m_1,m_2)\)。恢复条件 \(\gamma(m_1,m_2)=1\otimes(m_1,m_2)\) 因此等价于 \(m_2=f(m_1)\)，给出所述图子模。
\end{proof}

平坦而不忠实的映射则可能保留相容方程，却丢失原对象。例如 \(\mathbb Z\rightarrow\mathbb Q\) 把所有有限交换群张量成零；不同的原模于是具有相同的零下降数据，无法由它们唯一恢复。
